\documentclass[11pt]{article}
% ===========================================================================
%  Shared preamble for all MPM2D worksheets — input right after \documentclass.
%  (Files beginning with "_" are skipped by mpm2d-worksheet-publish.mjs.)
% ===========================================================================
\usepackage[letterpaper,top=2.2cm,bottom=1.9cm,left=1.6cm,right=1.6cm,headheight=28pt,footskip=24pt]{geometry}
\usepackage{amsmath,amssymb}
\usepackage{cancel}
\usepackage[T1]{fontenc}
\usepackage[utf8]{inputenc}
\usepackage{lmodern}
\usepackage{xcolor}
\usepackage[most]{tcolorbox}
\usepackage{enumitem}
\usepackage{multicol}
\usepackage{fancyhdr}
\usepackage{tikz}
\usetikzlibrary{calc}
\usepackage{pgfplots}
\pgfplotsset{compat=1.18}

\definecolor{exblue}{HTML}{4A90E2}
\definecolor{exbg}{HTML}{E6F3FF}
\definecolor{qorange}{HTML}{E69138}
\definecolor{qbg}{HTML}{FFF7CC}
\definecolor{keygray}{HTML}{475569}
\definecolor{keybg}{HTML}{F1F5F9}
\definecolor{iagreen}{HTML}{14653B}
\definecolor{titleblue}{HTML}{1D4ED8}
% Rotating palette so each worked example gets its own colour.
\definecolor{ex0}{HTML}{2563A0}\definecolor{ex1}{HTML}{3B7D3B}\definecolor{ex2}{HTML}{8A5A00}
\definecolor{ex3}{HTML}{A3327A}\definecolor{ex4}{HTML}{6D28A3}\definecolor{ex5}{HTML}{B08900}
\definecolor{ex6}{HTML}{0E7490}\definecolor{ex7}{HTML}{9A3412}\definecolor{ex8}{HTML}{15803D}

\pagestyle{fancy}
\fancyhf{}
\fancyhead[L]{\footnotesize NAME: \underline{\hspace{3.4cm}}}
\fancyhead[C]{\textbf{Integration Academy}\\[-2pt]{\scriptsize Math Simplified \textperiodcentered{} Success Amplified}}
\fancyhead[R]{\footnotesize MPM2D}
\fancyfoot[L]{\scriptsize dr.merie@IntegrationAcademy.ca}
\fancyfoot[C]{\scriptsize\thepage}
\fancyfoot[R]{\scriptsize IntegrationAcademy.ca}
\renewcommand{\headrulewidth}{1.2pt}
\renewcommand{\footrulewidth}{0.4pt}
\setlength{\parindent}{0pt}
\setlength{\parskip}{4pt}

% Examples are NOT breakable, so each example + its solution stays on one page.
\newtcolorbox{example}[2]{colframe=#1,colback=#1!7,coltitle=white,colbacktitle=#1,fonttitle=\bfseries,title={#2},arc=3pt,boxrule=1pt}
\newtcolorbox{qbox}[1]{colframe=qorange,colback=qbg,coltitle=white,colbacktitle=qorange,fonttitle=\bfseries,title={#1},arc=3pt,breakable,boxrule=1pt}
\newtcolorbox{keybox}{colframe=keygray,colback=keybg,coltitle=white,colbacktitle=keygray,fonttitle=\bfseries,title={$\checkmark$\ Answer Key},arc=3pt,breakable,boxrule=1pt}
\newtcolorbox{recapbox}{colframe=keygray,colback=keybg,coltitle=white,colbacktitle=keygray,fonttitle=\bfseries,title={Question Recap},arc=3pt,breakable,boxrule=1pt}
\newtcolorbox{ideabox}{colframe=titleblue,colback=titleblue!6,coltitle=white,colbacktitle=titleblue,fonttitle=\bfseries,title={The Big Ideas},arc=3pt,breakable,boxrule=1.2pt}
\newtcolorbox{learnbox}{colframe=iagreen,colback=iagreen!5,coltitle=white,colbacktitle=iagreen,fonttitle=\bfseries,title={The Big Picture},arc=3pt,breakable,boxrule=1.2pt}
\newcommand{\lh}[1]{\par\smallskip{\bfseries\color{iagreen}#1}\par\nobreak\smallskip}

\newcommand{\soln}{\par\smallskip\textit{\textbf{Solution.}}\ }
\newcommand{\workspace}[1]{\par\vspace{3pt}\fcolorbox{black!30}{white}{\begin{minipage}[t][#1][t]{\dimexpr\linewidth-2\fboxsep\relax}\ \end{minipage}}\par}
\newcommand{\sech}[1]{\par\vspace{8pt}{\Large\bfseries\color{iagreen}#1}\par\nobreak\vspace{2pt}{\color{black!20}\hrule height1pt}\vspace{6pt}}

% A compact coordinate plot sized for an example box: \eplot{xmin}{xmax}{ymin}{ymax}{body}
\newcommand{\eplot}[5]{\begin{center}\begin{tikzpicture}
\begin{axis}[width=6.8cm,height=4.4cm,axis lines=middle,xlabel={\tiny$x$},ylabel={\tiny$y$},
grid=both,grid style={gray!16},xmin=#1,xmax=#2,ymin=#3,ymax=#4,
xtick distance=2,ytick distance=2,tick label style={font=\tiny},axis line style={-},samples=120]
#5
\end{axis}\end{tikzpicture}\end{center}}

% A sine-curve plot (degrees on x, 0..360): \splot{ymin}{ymax}{body}
\newcommand{\splot}[3]{\begin{center}\begin{tikzpicture}
\begin{axis}[width=8.6cm,height=4.6cm,axis lines=middle,xlabel={\tiny$x\,(^\circ)$},ylabel={\tiny$y$},
grid=both,grid style={gray!16},xmin=0,xmax=360,ymin=#1,ymax=#2,
xtick distance=90,ytick distance=1,tick label style={font=\tiny},axis line style={-},samples=180]
#3
\end{axis}\end{tikzpicture}\end{center}}

% A small coordinate plot: \plot{xmin}{xmax}{ymin}{ymax}{body}
\newcommand{\plot}[5]{\begin{center}\begin{tikzpicture}
\begin{axis}[width=8.4cm,height=6cm,axis lines=middle,xlabel={\scriptsize$x$},ylabel={\scriptsize$y$},
grid=both,grid style={gray!16},xmin=#1,xmax=#2,ymin=#3,ymax=#4,
xtick distance=2,ytick distance=2,tick label style={font=\tiny},axis line style={-}]
#5
\end{axis}\end{tikzpicture}\end{center}}

% Right triangle (angle theta at bottom-left, right angle at bottom-right):
%   \rtri{drawBase}{drawHeight}{oppLabel}{adjLabel}{hypLabel}
\newcommand{\rtri}[5]{\begin{center}\begin{tikzpicture}[scale=0.85]
\coordinate (A) at (0,0);\coordinate (B) at (#1,0);\coordinate (C) at (#1,#2);
\draw[thick,exblue] (A)--(B)--(C)--cycle;
\draw[black] ($(B)+(-0.32,0)$)--($(B)+(-0.32,0.32)$)--($(B)+(0,0.32)$);
\node[below] at ($(A)!0.5!(B)$){#4};
\node[right] at ($(B)!0.5!(C)$){#3};
\node[above left] at ($(A)!0.5!(C)$){#5};
\node at (0.7,0.22){$\theta$};
\end{tikzpicture}\end{center}}

% General (oblique) triangle with vertex labels and opposite-side labels:
%   \gtri{Alabel}{Blabel}{Clabel}{aSide}{bSide}{cSide}
\newcommand{\gtri}[6]{\begin{center}\begin{tikzpicture}[scale=0.85]
\coordinate (A) at (0,0);\coordinate (B) at (5,0);\coordinate (C) at (1.6,3);
\draw[thick,exblue] (A)--(B)--(C)--cycle;
\node[below left] at (A){#1};\node[below right] at (B){#2};\node[above] at (C){#3};
\node[below] at ($(A)!0.5!(B)$){#6};
\node[right] at ($(B)!0.5!(C)$){#4};
\node[left] at ($(A)!0.5!(C)$){#5};
\end{tikzpicture}\end{center}}

\fancyhead[R]{\footnotesize MCV4U \textperiodcentered{} 6.5}
\begin{document}

\begin{center}
{\LARGE\bfseries\color{titleblue}Worksheet 6.5 --- The Cross Product}\\[2pt]
{\small Grade 12 Calculus and Vectors (MCV4U) \textperiodcentered{} Unit 6: Geometry and Algebra of Vectors}
\end{center}

\begin{ideabox}
$\vec a\times\vec b=(a_2b_3-a_3b_2,\,a_3b_1-a_1b_3,\,a_1b_2-a_2b_1)$ is perpendicular to both; $|\vec a\times\vec b|$ is the parallelogram area.
\begin{itemize}[leftmargin=*,itemsep=2pt]
  \item $\vec a\times\vec b=(a_2b_3-a_3b_2,\,a_3b_1-a_1b_3,\,a_1b_2-a_2b_1)$.
  \item $\vec a\times\vec b\perp\vec a$ and $\perp\vec b$.
  \item $|\vec a\times\vec b|=$ parallelogram area; triangle $=\tfrac12|\vec a\times\vec b|$.
\end{itemize}
\end{ideabox}

\begin{learnbox}
\lh{Multiplying to a vector}The \textbf{cross product} $\vec u\times\vec v$ produces a \emph{vector} perpendicular to both $\vec u$ and $\vec v$ in 3-D.
\lh{Magnitude and area}Its magnitude $|\vec u|\,|\vec v|\sin\theta$ equals the area of the parallelogram the two vectors span.
\lh{Direction}The right-hand rule fixes the direction, and $\vec u\times\vec v=-\,\vec v\times\vec u$ — order matters.
\end{learnbox}

\sech{Worked Examples}

\begin{example}{ex0}{Example 1 --- Standard basis}
Compute $(1,0,0)\times(0,1,0)$.\soln
\textbf{Step 1 (apply the cross-product formula $(a_2b_3-a_3b_2,\,a_3b_1-a_1b_3,\,a_1b_2-a_2b_1)$ with $\vec a=(1,0,0),\vec b=(0,1,0)$):} first component $0\cdot0-0\cdot1=0$; second $0\cdot0-1\cdot0=0$; third $1\cdot1-0\cdot0=1$.\\[3pt]
\textbf{Conclusion:} $(1,0,0)\times(0,1,0)=(0,0,1)$.
\end{example}

\begin{example}{ex1}{Example 2 --- General}
Compute $(1,2,3)\times(4,5,6)$.\soln
\textbf{Step 1 (label components):} $\vec a=(1,2,3)$, $\vec b=(4,5,6)$.\\[3pt]
\textbf{Step 2 (apply the formula component by component):} first: $2\cdot6-3\cdot5=12-15=-3$; second: $3\cdot4-1\cdot6=12-6=6$; third: $1\cdot5-2\cdot4=5-8=-3$.\\[3pt]
\textbf{Conclusion:} $(1,2,3)\times(4,5,6)=(-3,6,-3)$.
\end{example}

\begin{example}{ex2}{Example 3 --- Perpendicular check}
Verify $(-3,6,-3)\perp(1,2,3)$.\soln
\textbf{Step 1 (use the dot product to test perpendicularity):} $(-3)(1)+(6)(2)+(-3)(3)$.\\[3pt]
\textbf{Step 2 (compute):} $-3+12-9$.\\[3pt]
\textbf{Conclusion:} $=0$, so the vectors are perpendicular (as the cross product must be, to both original vectors).
\end{example}

\begin{example}{ex3}{Example 4 --- Parallelogram area}
Area from $(1,0,0)$ and $(0,2,0)$.\soln
\textbf{Step 1 (compute the cross product):} $(1,0,0)\times(0,2,0)=(0\cdot0-0\cdot2,\,0\cdot0-1\cdot0,\,1\cdot2-0\cdot0)=(0,0,2)$.\\[3pt]
\textbf{Step 2 (the parallelogram area equals the magnitude of the cross product):} $|(0,0,2)|=2$.\\[3pt]
\textbf{Conclusion:} area $=2$.
\end{example}

\begin{example}{ex4}{Example 5 --- Triangle area}
Triangle area from $(1,0,0)$ and $(0,2,0)$.\soln
\textbf{Step 1 (recall the triangle is half the parallelogram):} from the previous example, the parallelogram area is $2$.\\[3pt]
\textbf{Step 2 (halve it):} $\tfrac12(2)$.\\[3pt]
\textbf{Conclusion:} triangle area $=1$.
\end{example}

\begin{example}{ex5}{Example 6 --- Standard basis}
Compute $(0,1,0)\times(0,0,1)$.\soln
\textbf{Step 1 (apply the formula with $\vec a=(0,1,0),\vec b=(0,0,1)$):} first: $1\cdot1-0\cdot0=1$; second: $0\cdot0-0\cdot1=0$; third: $0\cdot0-1\cdot0=0$.\\[3pt]
\textbf{Conclusion:} $(0,1,0)\times(0,0,1)=(1,0,0)$.
\end{example}

\begin{example}{ex6}{Example 7 --- General}
Compute $(2,0,0)\times(0,3,0)$.\soln
\textbf{Step 1 (apply the formula):} first: $0\cdot0-0\cdot3=0$; second: $0\cdot0-2\cdot0=0$; third: $2\cdot3-0\cdot0=6$.\\[3pt]
\textbf{Conclusion:} $(2,0,0)\times(0,3,0)=(0,0,6)$.
\end{example}

\begin{example}{ex7}{Example 8 --- General}
Compute $(1,1,0)\times(0,1,1)$.\soln
\textbf{Step 1 (apply the formula with $\vec a=(1,1,0),\vec b=(0,1,1)$):} first: $1\cdot1-0\cdot1=1$; second: $0\cdot0-1\cdot1=-1$; third: $1\cdot1-1\cdot0=1$.\\[3pt]
\textbf{Conclusion:} $(1,1,0)\times(0,1,1)=(1,-1,1)$.
\end{example}

\begin{example}{ex8}{Example 9 --- Parallelogram area}
Area from $(3,0,0)$ and $(0,4,0)$.\soln
\textbf{Step 1 (compute the cross product):} $(3,0,0)\times(0,4,0)=(0\cdot0-0\cdot4,\,0\cdot0-3\cdot0,\,3\cdot4-0\cdot0)=(0,0,12)$.\\[3pt]
\textbf{Step 2 (take the magnitude):} $|(0,0,12)|=12$.\\[3pt]
\textbf{Conclusion:} area $=12$.
\end{example}

\clearpage
\sech{Your Turn --- Show Your Work}
For each question, show all steps clearly in the space provided.

\begin{qbox}{Question 1}
Compute $(0,1,0)\times(0,0,1)$.\workspace{2.4cm}
\end{qbox}

\begin{qbox}{Question 2}
Compute $(2,0,0)\times(0,3,0)$.\workspace{2.4cm}
\end{qbox}

\begin{qbox}{Question 3}
What direction is $\vec a\times\vec b$ relative to $\vec a,\vec b$?\workspace{2.4cm}
\end{qbox}

\begin{qbox}{Question 4}
Area of the parallelogram from $(3,0,0)$ and $(0,4,0)$?\workspace{2.4cm}
\end{qbox}

\begin{qbox}{Question 5}
Compute $(1,1,0)\times(0,1,1)$.\workspace{2.4cm}
\end{qbox}

\begin{qbox}{Question 6}
Compute $(1,0,0)\times(0,0,1)$.\workspace{2.4cm}
\end{qbox}

\begin{qbox}{Question 7}
Compute $(1,2,3)\times(4,5,6)$.\workspace{2.4cm}
\end{qbox}

\begin{qbox}{Question 8}
Triangle area from $(2,0,0)$ and $(0,2,0)$?\workspace{2.4cm}
\end{qbox}

\begin{qbox}{Question 9}
Is $(1,2,3)\times(2,4,6)$ the zero vector? Why?\workspace{2.4cm}
\end{qbox}

\begin{qbox}{Question 10}
Compute $(0,0,1)\times(1,0,0)$.\workspace{2.4cm}
\end{qbox}

\begin{qbox}{Question 11}
Magnitude of $(1,0,0)\times(0,5,0)$?\workspace{2.4cm}
\end{qbox}

\begin{qbox}{Question 12}
Verify $(1,2,3)\times(4,5,6)=(-3,6,-3)$ is perpendicular to $(4,5,6)$.\workspace{2.4cm}
\end{qbox}

\begin{qbox}{Question 13 --- Challenge}
Find the area of the triangle with vertices $A(0,0,0)$, $B(1,0,0)$, $C(0,2,3)$.\workspace{3.5cm}
\end{qbox}

\clearpage
\sech{Question Recap \& Answer Key}

\begin{recapbox}
\begin{multicols}{2}
\small
\begin{enumerate}[leftmargin=*,itemsep=3pt]
  \item Compute $(0,1,0)\times(0,0,1)$.
  \item Compute $(2,0,0)\times(0,3,0)$.
  \item What direction is $\vec a\times\vec b$ relative to $\vec a,\vec b$?
  \item Area of the parallelogram from $(3,0,0)$ and $(0,4,0)$?
  \item Compute $(1,1,0)\times(0,1,1)$.
  \item Compute $(1,0,0)\times(0,0,1)$.
  \item Compute $(1,2,3)\times(4,5,6)$.
  \item Triangle area from $(2,0,0)$ and $(0,2,0)$?
  \item Is $(1,2,3)\times(2,4,6)$ the zero vector? Why?
  \item Compute $(0,0,1)\times(1,0,0)$.
  \item Magnitude of $(1,0,0)\times(0,5,0)$?
  \item Verify $(1,2,3)\times(4,5,6)=(-3,6,-3)$ is perpendicular to $(4,5,6)$.
  \item Find the area of the triangle with vertices $A(0,0,0)$, $B(1,0,0)$, $C(0,2,3)$.
\end{enumerate}
\end{multicols}
\end{recapbox}

\begin{keybox}
\begin{multicols}{2}
\small
\begin{enumerate}[leftmargin=*,itemsep=3pt]
  \item $(1,0,0)$
  \item $(0,0,6)$
  \item perpendicular to both
  \item $12$
  \item $(1,-1,1)$
  \item $(0,-1,0)$
  \item $(-3,6,-3)$
  \item $\tfrac12(4)=2$
  \item yes — they are parallel ($\sin\theta=0$)
  \item $(0,1,0)$
  \item $5$
  \item $-12+30-18=0$ ✓
  \item $\vec{AB}\times\vec{AC}=(1,0,0)\times(0,2,3)=(0,-3,2)$; area $=\tfrac12\sqrt{13}$
\end{enumerate}
\end{multicols}
\end{keybox}

\end{document}
